\documentclass[11pt,a4paper]{article}

\usepackage[utf8]{inputenc}
\usepackage[T1]{fontenc}
\usepackage{amsmath,amssymb}
\usepackage[italian]{babel}
\usepackage[margin=2.3cm]{geometry}
\usepackage{xcolor}
\usepackage{listings}
\usepackage{booktabs}
\usepackage{tabularx}
\usepackage{longtable}
\usepackage{enumitem}
\usepackage{fancyhdr}
\usepackage{titlesec}
\usepackage{pifont}
\usepackage[hidelinks]{hyperref}

\hypersetup{colorlinks=true, linkcolor=blue!55!black, urlcolor=blue!55!black}

% ---- colori & stile codice ----
\definecolor{codebg}{RGB}{245,246,248}
\definecolor{codekw}{RGB}{0,90,160}
\definecolor{codecm}{RGB}{100,120,100}
\definecolor{codest}{RGB}{160,60,40}
\definecolor{accent}{RGB}{20,90,150}

\lstdefinestyle{cpp}{
  language=C++, backgroundcolor=\color{codebg}, basicstyle=\ttfamily\small,
  keywordstyle=\color{codekw}\bfseries, commentstyle=\color{codecm}\itshape,
  stringstyle=\color{codest}, breaklines=true, showstringspaces=false,
  frame=leftline, framerule=1.2pt, rulecolor=\color{accent},
  xleftmargin=8pt, aboveskip=8pt, belowskip=6pt,
  morekeywords={Ptr,Chunk,Packet,uint64_t,override,virtual,simtime_t}
}
\lstdefinestyle{sh}{
  backgroundcolor=\color{codebg}, basicstyle=\ttfamily\footnotesize,
  breaklines=true, frame=leftline, framerule=1.2pt, rulecolor=\color{codecm},
  xleftmargin=8pt
}
\lstset{style=cpp}

\newcommand{\code}[1]{\texttt{#1}}

\titleformat{\section}{\Large\bfseries\color{accent}}{\thesection}{0.6em}{}
\titleformat{\subsection}{\large\bfseries}{\thesubsection}{0.6em}{}

\pagestyle{fancy}\fancyhf{}
\lhead{\small QUIC DATAGRAM per TOD-over-5G}
\rhead{\small\thepage}
\renewcommand{\headrulewidth}{0.4pt}

\title{\vspace{-1cm}\textbf{Implementazione di QUIC DATAGRAM (RFC 9221) in INET}\\[2mm]
\large per il trasporto inaffidabile dei dati-sensore nel simulatore TOD-over-5G\\(CARLA + OMNeT++/Simu5G)}
\author{Progetto TOD --- Simulazione Tele-Operated Driving su 5G}
\date{\today}

\begin{document}
\maketitle
\thispagestyle{fancy}

\begin{abstract}
\noindent Questo documento descrive in modo completo l'estensione del modulo di trasporto
\textbf{QUIC di INET~4.6.0} con i \textbf{DATAGRAM frame} definiti dalla RFC~9221, realizzata per
sostituire gli stream QUIC affidabili con un canale \emph{inaffidabile e non ordinato} nel trasporto
dei dati-sensore del veicolo tele-operato. Vengono documentate tutte le classi e tutti i metodi
introdotti o modificati, il meccanismo che garantisce l'assenza di ritrasmissione, la negoziazione del
parametro \code{max\_datagram\_frame\_size}, le applicazioni di test, i risultati della verifica a runtime
e la terminologia di rete usata.
\end{abstract}

\tableofcontents
\newpage

% =====================================================================
\section{Contesto e motivazione}
Il simulatore TOD (Tele-Operated Driving) accoppia \textbf{CARLA} (mondo fisico e percezione) e
\textbf{OMNeT++/Simu5G} (rete 5G) tramite la libreria CARLANeT. Il veicolo campiona i sensori (camera,
LiDAR, radar) e li invia, attraverso la rete 5G, all'agente tele-operatore remoto, che calcola il comando
di guida. La metrica chiave \`e il \code{loss\_ratio}: la frazione di dati-sensore persi per ogni
\emph{sensor frame} (un campionamento a un dato istante).

\medskip
\noindent\textbf{Il problema.} Nella versione precedente ogni sensore inviava i propri dati su uno
\textbf{stream QUIC dedicato e persistente}. Poich\'e QUIC garantisce consegna \emph{affidabile e ordinata}
all'interno di ogni stream, la perdita di un singolo pacchetto \textbf{blocca} (head-of-line blocking) tutti
i frame successivi di quel sensore finch\'e non arriva la ritrasmissione. Di conseguenza una singola
perdita si propaga a pi\`u \emph{sensor frame} consecutivi, gonfiando artificialmente il \code{loss\_ratio}
con perdite correlate. Per uno streaming di sensori in tempo reale ci\`o non \`e realistico: un frame
ritrasmesso arriva ormai scaduto e inutile, ma intanto ha bloccato quelli nuovi.

\medskip
\noindent\textbf{La soluzione.} Introdurre i \textbf{QUIC DATAGRAM frame} (RFC~9221): un servizio di
consegna \emph{inaffidabile}, \emph{non ordinato} e \emph{senza ritrasmissione}, ideale per i dati-sensore.
Il canale di controllo/stato resta su uno stream affidabile (stream~0). Poich\'e il QUIC di INET non
implementava i DATAGRAM, sono stati aggiunti al modulo di trasporto.

% =====================================================================
\section{Terminologia}
Le parole ``datagram'', ``frame'' e ``packet'' sono sovraccariche: significano cose diverse a livelli
diversi della pila. La Tabella~\ref{tab:stack} e le successive chiariscono l'uso nel progetto.

\begin{table}[h]\centering\small
\caption{La pila protocollare: una unit\`a per livello.}\label{tab:stack}
\begin{tabularx}{\textwidth}{@{}l l X@{}}
\toprule
\textbf{Livello} & \textbf{Unit\`a} & \textbf{Nella nostra implementazione}\\
\midrule
App (\code{TODCarApp}) & messaggio / \emph{sensor frame} & \code{SensorDataResponse} per sensore; \code{frameId} = batch di un tick\\
QUIC (transport) & \textbf{QUIC frame} & l'unit\`a con i dati: STREAM, DATAGRAM, ACK, MAX\_DATA\dots\\
 & \textbf{QUIC packet} & bundle di frame, dimensionato $\le$ PMTU, con packet number\\
UDP & UDP datagram & uno per ogni QUIC packet\\
IP & IP packet (= IP datagram) & con bit DF $\Rightarrow$ non frammentato\\
Link (PPP/Eth sim) & link-layer frame & PPP frame\\
\bottomrule
\end{tabularx}
\end{table}

\begin{table}[h]\centering\small
\caption{La parola ``datagram'' a tre livelli diversi.}
\begin{tabularx}{\textwidth}{@{}l l X@{}}
\toprule
\textbf{Significa} & \textbf{Livello} & \textbf{Nel progetto}\\
\midrule
\textbf{QUIC DATAGRAM frame} (RFC 9221) & QUIC & \ding{51}~\textbf{QUESTO} --- la feature implementata\\
UDP datagram & UDP & il contenitore di un QUIC packet\\
IP datagram & IP & sinonimo di IP packet\\
\bottomrule
\end{tabularx}
\end{table}

\begin{table}[h]\centering\small
\caption{La parola ``frame'', usata in due sensi diversi.}
\begin{tabularx}{\textwidth}{@{}l l X@{}}
\toprule
\textbf{Significa} & \textbf{Livello} & \textbf{Nel progetto}\\
\midrule
\textbf{QUIC frame} & QUIC & l'unit\`a dentro un QUIC packet (STREAM, DATAGRAM, ACK\dots)\\
\textbf{sensor frame} (\code{frameId}) & App & batch di dati-sensore di un tick; base del \code{loss\_ratio}\\
link-layer frame & Data-link & PPP/Ethernet nel simulatore\\
\bottomrule
\end{tabularx}
\end{table}

\noindent\textbf{Spezzare i dati: tre meccanismi diversi.}
\begin{itemize}[nosep,leftmargin=1.4em]
  \item \textbf{Segmentazione} (stream) --- la fa \emph{QUIC}, automaticamente: un messaggio stream $\to$
    pi\`u STREAM frame su pi\`u QUIC packet, riassemblati al ricevitore.
  \item \textbf{Frammentazione IP} --- la farebbe \emph{IP} su un pacchetto oversized: QUIC la \textbf{evita}
    (bit DF + dimensionamento a PMTU). Non accade.
  \item \textbf{Chunking} (datagram) --- lo fa l'\emph{applicazione}, manualmente: QUIC non frammenta un
    datagram, quindi \`e l'app a spezzare il messaggio in pezzi $\le$ \code{max\_datagram\_frame\_size}.
\end{itemize}

\noindent\textbf{Glossario di progetto.}
\begin{description}[leftmargin=0pt,style=nextline,itemsep=1pt]
  \item[\code{loss\_ratio}] frazione di dati-sensore persi per un \emph{sensor frame}: $1-\text{arrivati}/\text{attesi}$. Calcolato da \code{TODAgentApp}, inviato a CARLA, usato per la soglia $\to$ arresto d'emergenza.
  \item[stream 0] lo \emph{stream QUIC} affidabile e ordinato che porta status/controllo.
  \item[HoL blocking] head-of-line blocking: in uno stream, un pacchetto perso blocca la consegna dei successivi finch\'e non \`e ritrasmesso. I datagram lo evitano.
  \item[\code{max\_datagram\_frame\_size}] transport parameter RFC 9221: dimensione massima di un DATAGRAM frame che il peer accetta (0 = non supportati).
  \item[PMTU / DPLPMTUD] Path MTU e il meccanismo con cui QUIC la scopre (RFC 8899) per non far frammentare IP.
\end{description}

% =====================================================================
\section{Architettura: il percorso end-to-end}
Il datagram attraversa i seguenti passaggi, dall'invio applicativo fino alla consegna al ricevitore:

\begin{lstlisting}[style=sh]
App:  QuicSocket::sendDatagram(pkt)                     [contract/quic/QuicSocket]
        | setKind(QUIC_C_SEND_DATAGRAM)
        v
ConnectionState::processFrame? no -> processAppCommand
   EstablishedConnectionState::processSendDatagramAppCommand
        v
Connection::newDatagram(data)                           [quic/Connection]
        | verifica max_datagram_frame_size del peer
        | crea DatagramFrameHeader(length) + QuicDatagramFrame -> datagramQueue
        v
PacketBuilder::addFramesFromDatagramQueue (in buildPacket)   [quic/PacketBuilder]
        | size-aware; scarta se non entra nel pacchetto (no frammentazione)
        v ...UDP / IP / rete 5G lossy...  (frame perso => onFrameLost() no-op => NIENTE retx)
RX: ConnectionState::processFrame -> case FRAME_HEADER_TYPE_DATAGRAM
        v
EstablishedConnectionState::processDatagramFrame
        | popAtFront(B(length))  <- legge esattamente il payload
        v
Connection::processReceivedDatagram -> AppSocket::sendDatagram   [quic/AppSocket]
        | QUIC_I_DATAGRAM
        v
App:  QuicSocket::ICallback::socketDatagramArrived()
\end{lstlisting}

% =====================================================================
\section{Modifiche a INET: classi e metodi}
Tutte le modifiche sono nel modulo \code{src/inet/transportlayer/quic} (e nel contract
\code{transportlayer/contract/quic}). Sono versionate come patch in \code{tod-omnet/inet-patches/}.

\subsection{Nuovo tipo di frame (\code{packet/FrameHeader.msg})}
\begin{itemize}[nosep,leftmargin=1.4em]
  \item Nuovo valore di enum \code{FRAME\_HEADER\_TYPE\_DATAGRAM = 0x31} (la variante RFC 9221 ``con Length'').
  \item Nuova classe chunk \code{DatagramFrameHeader} (estende \code{FrameHeader}): un solo campo
    \code{VariableLengthInteger length} (lunghezza in byte del payload). A differenza dello
    \code{StreamFrameHeader} non ha \code{streamId} n\'e \code{offset}.
    \begin{itemize}[nosep]
      \item \code{calcChunkLength()} --- calcola la lunghezza dell'header: $1 + \text{varint}(\text{length})$ byte.
      \item \code{str()} --- rappresentazione testuale \code{DATAGRAM[length=\dots]} per il logging.
    \end{itemize}
  \item Alla classe \code{TransportParametersExtension} \`e stato aggiunto il campo
    \code{VariableLengthInteger maxDatagramFrameSize} (con aggiornamento di \code{calcChunkLength()}),
    per la negoziazione RFC 9221 durante l'handshake.
\end{itemize}

\subsection{\code{QuicDatagramFrame} (\code{packet/QuicDatagramFrame.\{h,cc\}}) --- nuova}
Rappresenta un DATAGRAM frame. A differenza di \code{QuicStreamFrame} (che recupera i dati ``lazy'' da uno
\code{Stream}) tiene il payload \emph{direttamente} come data del frame. Metodi:
\begin{itemize}[nosep,leftmargin=1.4em]
  \item \code{QuicDatagramFrame()} --- costruttore di default.
  \item \code{QuicDatagramFrame(Ptr<const DatagramFrameHeader> header, Ptr<const Chunk> data)} ---
    costruisce header+payload; usa chiamate esplicite a \code{setHeader}/\code{setData} (non delega al
    costruttore base) per far valorizzare correttamente \code{datagramHeader} tramite dispatch virtuale.
  \item \code{setHeader(Ptr<const FrameHeader>) override} --- invoca la base e memorizza l'header
    come \code{DatagramFrameHeader} (\code{dynamicPtrCast}).
  \item \code{hasData() override} $\to$ \code{true} --- garantisce che il payload venga serializzato nel packet.
  \item \code{equals(const QuicFrame*) override} --- confronto per \emph{identit\`a}: un datagram \`e unico e
    non ritrasmesso, non ha ``identit\`a di ritrasmissione''.
  \item \code{getDatagramHeader()} --- accessor all'header tipizzato.
  \item \textbf{Nessun override di \code{onFrameLost()}} --- \`e questo che disabilita la ritrasmissione
    (vedi \S\ref{sec:noretx}). Eredita inoltre \code{isAckEliciting()=true} e \code{countsAsInFlight()=true}
    dal comportamento di default di \code{QuicFrame::setHeader}, corretto per RFC 9221.
\end{itemize}

\subsection{\code{Connection} (\code{Connection.\{h,cc\}})}
\begin{itemize}[nosep,leftmargin=1.4em]
  \item Nuovo membro \code{std::vector<QuicFrame*> datagramQueue} --- coda dei DATAGRAM frame in attesa di invio.
  \item \code{getDatagramQueue()} --- accessor alla coda (usato dal \code{PacketBuilder}).
  \item Nel costruttore: \code{packetBuilder->setDatagramQueue(\&datagramQueue)} collega la coda al builder.
  \item \code{newDatagram(Ptr<const Chunk> data)} --- accoda un datagram per l'invio:
    \begin{enumerate}[nosep,label=\arabic*.]
      \item Se il peer non supporta i datagram (\code{maxDatagramFrameSize == 0}) $\to$ scarta (RFC 9221).
      \item Crea \code{DatagramFrameHeader} con \code{length} = byte del payload e \code{calcChunkLength()}.
      \item Crea il \code{QuicDatagramFrame}; se la sua dimensione (header+payload) eccede il
        \code{max\_datagram\_frame\_size} del peer $\to$ scarta.
      \item Se la coda ha gi\`a \code{maxDatagramQueueLength} elementi $\to$ scarta il \textbf{pi\`u
        vecchio} (in testa) per fare posto, ed emette \code{datagramDropped}. Vedi
        \S\ref{sec:queuepolicy}.
      \item Accoda in \code{datagramQueue} e chiama \code{sendPackets()}; poi emette
        \code{datagramQueueLength} (il backlog rimasto dopo il tentativo di invio).
    \end{enumerate}
  \item \code{processReceivedDatagram(Ptr<const Chunk> data)} --- consegna immediata all'app tramite
    \code{appSocket->sendDatagram(data)} (nessun reassembly, nessun ordinamento).
\end{itemize}

\subsection{\code{PacketBuilder} (\code{PacketBuilder.\{h,cc\}})}
\begin{itemize}[nosep,leftmargin=1.4em]
  \item Nuovo membro \code{datagramQueue} e setter \code{setDatagramQueue(...)}.
  \item \code{addFramesFromDatagramQueue(QuicPacket* packet, int maxPacketSize)} --- drena la coda dei
    datagram nel pacchetto in costruzione, in modo \emph{size-aware} usando \code{frame->getSize()}
    (header+payload, a differenza dei control frame che sono solo header). Politica RFC 9221:
    \begin{itemize}[nosep]
      \item se un datagram non entra nemmeno in un pacchetto vuoto $\to$ viene \textbf{scartato}
        (non si frammenta);
      \item se non entra nel pacchetto \emph{corrente} (gi\`a parzialmente pieno) $\to$ resta in coda per il
        pacchetto successivo;
      \item altrimenti viene aggiunto al pacchetto.
    \end{itemize}
  \item \code{buildPacket(...)} --- chiama \code{addFramesFromDatagramQueue} \textbf{per ultimo}, dopo
    control frame e frame STREAM, cos\`i i datagram riempiono lo spazio che avanza nei normali
    pacchetti dati. Vedi \S\ref{sec:queuepolicy}.
\end{itemize}

\subsection{Politica di invio: priorit\`a e coda}
\label{sec:queuepolicy}

I DATAGRAM sono traffico bulk best-effort con una scadenza; gli STREAM portano il canale
affidabile. Due scelte, indipendenti dalla RFC ma necessarie perch\'e una sorgente datagram satura
non distrugga la connessione:

\begin{description}[nosep,leftmargin=1.4em]
  \item[Gli STREAM hanno priorit\`a stretta sui DATAGRAM.] In \code{buildPacket} i DATAGRAM frame si
    impacchettano per ultimi. Nell'ordine opposto un frammento datagram da $\sim$1\,kB riempie da solo
    il pacchetto e non lascia mai spazio a un frame STREAM: lo stream affidabile viene affamato a
    tempo indeterminato. In un sistema di teleguida significa perdere il canale di controllo
    (status e istruzioni su stream 0) proprio quando la rete \`e congestionata, cio\`e quando serve.
    Il rovescio \`e che ora \`e uno stream saturo a poter affamare i datagram: \`e la politica voluta
    (\emph{control plane} sopra \emph{sensor plane}).

  \item[La coda datagram \`e limitata e scarta i pi\`u vecchi.] Parametro NED
    \code{maxDatagramQueueLength} (default 64; 0 = illimitata). Un DATAGRAM non \`e ritrasmesso e
    trasporta dati con una scadenza: con una coda illimitata, un'app che offre pi\`u di quanto il path
    riesca a portare accumula un backlog che non drena mai e finisce per spedire \emph{solo} dati
    obsoleti, che il ricevitore scarta come fuori deadline. Scartando il pi\`u vecchio, sul link
    viaggia sempre il dato pi\`u fresco.
\end{description}

Misura su \code{examples/quic/datagram\_test}, \code{[Config Overload]} (sorgente continua a
$17\times$ la capacit\`a del link, nessuna perdita di canale, buffer del router 20 pacchetti):
a parit\`a di datagram consegnati (291 in 10\,s, cio\`e la capacit\`a del link), con coda illimitata
l'ultimo dato consegnato \`e vecchio di \textbf{9{,}3\,s}; con la coda limitata a 8, di
\textbf{0{,}56\,s}.

Statistiche per connessione: \code{datagramQueueLength} (backlog dopo ogni tentativo di invio) e
\code{datagramDropped} (datagram scartati per coda piena).

\subsection{\code{ConnectionState} / \code{EstablishedConnectionState} (\code{connectionstate/})}
La macchina a stati fa il dispatch dei comandi applicativi e dei frame ricevuti.
\begin{itemize}[nosep,leftmargin=1.4em]
  \item In \code{ConnectionState::processAppCommand} (switch): nuovo case \code{QUIC\_C\_SEND\_DATAGRAM} $\to$
    \code{processSendDatagramAppCommand(msg)}.
  \item In \code{ConnectionState::processFrame} (switch sul tipo di frame ricevuto): nuovo case
    \code{FRAME\_HEADER\_TYPE\_DATAGRAM} $\to$ \code{processDatagramFrame(...)}.
  \item \code{ConnectionState::processSendDatagramAppCommand(cMessage*)} (base, \code{virtual}) --- lancia
    ``unexpected in the current state'' (i datagram si inviano solo da stato Established).
  \item \code{ConnectionState::processDatagramFrame(const Ptr<const DatagramFrameHeader>\&, Packet*)}
    (base, \code{virtual}) --- idem.
  \item \code{EstablishedConnectionState::processSendDatagramAppCommand(cMessage*) override} --- estrae il
    \code{Packet}, chiama \code{context->newDatagram(pkt->peekData())}.
  \item \code{EstablishedConnectionState::processDatagramFrame(header, pkt) override} --- legge
    \emph{esattamente} \code{length} byte con \code{pkt->popAtFront(B(header->getLength()))} (un DATAGRAM
    frame pu\`o essere seguito da altri frame nello stesso pacchetto) e chiama
    \code{context->processReceivedDatagram(data)}.
\end{itemize}

\subsection{\code{AppSocket} (\code{AppSocket.\{h,cc\}})}
\begin{itemize}[nosep,leftmargin=1.4em]
  \item \code{sendDatagram(Ptr<const Chunk> data)} --- costruisce un \code{Packet} con kind
    \code{QUIC\_I\_DATAGRAM}, vi inserisce il payload e lo invia verso l'app (indicazione dedicata,
    parallela a \code{sendData}/\code{QUIC\_I\_DATA} usata per gli stream).
\end{itemize}

\subsection{\code{QuicSocket} + \code{QuicCommand.msg} (contract) --- API applicativa}
\begin{itemize}[nosep,leftmargin=1.4em]
  \item \code{QuicCommand.msg}: nuovo command code \code{QUIC\_C\_SEND\_DATAGRAM = 9} e nuova indicazione
    \code{QUIC\_I\_DATAGRAM = 12}.
  \item \code{QuicSocket::sendDatagram(Packet* msg)} --- imposta \code{setKind(QUIC\_C\_SEND\_DATAGRAM)} e invia
    al modulo QUIC (nessun tag di stream: il comando dedicato instrada verso \code{Connection::newDatagram}).
  \item \code{QuicSocket::ICallback::socketDatagramArrived(QuicSocket*, Packet*)} --- nuova callback,
    dichiarata \emph{non-pure} con default vuoto (\code{delete packet;}): le app QUIC esistenti non vanno
    modificate.
  \item \code{QuicSocket::processMessage} --- nuovo case \code{QUIC\_I\_DATAGRAM} che invoca
    \code{cb->socketDatagramArrived(this, packet)} (rispecchia \code{QUIC\_I\_DATA}).
\end{itemize}

\subsection{\code{TransportParameters} e \code{Quic.ned} --- negoziazione}
\begin{itemize}[nosep,leftmargin=1.4em]
  \item \code{TransportParameters}: nuovo membro \code{uint64\_t maxDatagramFrameSize}.
  \item \code{readParameters(Quic*)} --- legge il valore locale dal parametro NED
    \code{par("maxDatagramFrameSize")}.
  \item \code{readExtension(...)} --- legge il valore annunciato dal peer da
    \code{TransportParametersExtension::getMaxDatagramFrameSize()}.
  \item \code{Quic.ned}: nuovo parametro \code{int maxDatagramFrameSize @unit(B) = default(1200B)}
    (0 disabilita i datagram). Il valore locale viene serializzato nell'extension e inviato durante
    l'handshake da \code{PacketBuilder}; il valore del peer viene salvato in \code{remoteTransportParameters}
    e usato da \code{Connection::newDatagram} per l'enforcement.
\end{itemize}

% =====================================================================
\section{Il meccanismo di non-ritrasmissione}\label{sec:noretx}
La propriet\`a pi\`u importante dei datagram \`e che \textbf{un datagram perso non viene ritrasmesso}. In INET
questo si ottiene \emph{senza toccare il gestore di affidabilit\`a}, grazie al polimorfismo di
\code{QuicFrame}.

Quando un pacchetto viene dichiarato perso, per ogni frame contenuto viene chiamato
\code{frame->onFrameLost()} (in \code{QuicPacket.cc}). Le implementazioni concrete decidono cosa fare:
\begin{itemize}[nosep,leftmargin=1.4em]
  \item \code{QuicStreamFrame::onFrameLost()} ri-accoda i dati per la ritrasmissione;
  \item la base \code{QuicFrame::onFrameLost()} \`e un \textbf{no-op} (corpo vuoto).
\end{itemize}
Poich\'e \code{QuicDatagramFrame} \textbf{non} sovrascrive \code{onFrameLost()}, eredita il no-op: il payload
perso semplicemente \emph{sparisce}. In pi\`u, il percorso di ritrasmissione nei probe packet
(\code{PacketBuilder::buildAckElicitingPacket}) filtra esplicitamente solo i frame di tipo
\code{FRAME\_HEADER\_TYPE\_STREAM}, quindi i datagram sono esclusi anche l\`i. Il datagram resta comunque
\emph{ack-eliciting} (contribuisce a RTT e congestion control), come prescrive la RFC~9221.

% =====================================================================
\section{Applicazioni di test e rete}
Per la verifica funzionale sono state scritte due app minimali (in \code{applications/quicapp/}) e una rete
con perdita (in \code{examples/quic/datagram\_test/}).

\subsection{\code{QuicDatagramSender}}
Invia $N$ datagrammi numerati a intervalli fissi. Metodi principali:
\code{handleStartOperation} (bind + connect + crea il timer), \code{socketEstablished} (avvia l'invio dopo
l'handshake), \code{sendOneDatagram} (crea un \code{ApplicationPacket} con \code{sequenceNumber} e chiama
\code{socket.sendDatagram}), \code{finish} (registra lo scalare \code{datagramsSent}). Parametri:
\code{numDatagrams}, \code{datagramBytes}, \code{sendInterval}, \code{connectAddress/Port}.

\subsection{\code{QuicDatagramReceiver}}
Ascolta, accetta la connessione e conta i datagrammi ricevuti. Metodi principali:
\code{handleStartOperation} (bind + listen), \code{socketConnectionAvailable} (accept),
\code{socketDatagramArrived} (legge il \code{sequenceNumber}, logga eventuali \emph{gap} = perdite, conta),
\code{finish} (registra \code{datagramsReceived}).

\subsection{Rete \code{DatagramTest.ned}}
Topologia \code{sender -- router1 =lossy= router2 -- receiver}. Il link centrale \`e un
\code{PeriodicLossChannel} con parametro \code{per} (packet error rate) applicato in direzione
sender$\to$receiver, cos\`i gli ACK di ritorno non vengono persi e la perdita colpisce solo i datagram.

% =====================================================================
\section{Risultati della verifica}
Esecuzione con \textbf{20\% di perdita} sul link centrale, 20 datagrammi da 800\,B:

\begin{table}[h]\centering\small
\begin{tabular}{@{}ll@{}}
\toprule
Datagrammi inviati & 20\\
Datagrammi ricevuti & 16\\
Persi (mai ritrasmessi) & seq \textbf{2, 7, 13, 19}\\
Dopo ogni ``buco'' & il seq successivo arriva \emph{subito}\\
\bottomrule
\end{tabular}
\end{table}

\noindent Sequenza ricevuta: \code{0,1,\_,3,4,5,6,\_,8,9,10,11,12,\_,14,15,16,17,18}. Questo dimostra
sperimentalmente tutte le propriet\`a attese:
\begin{itemize}[nosep,leftmargin=1.4em]
  \item \textbf{Consegna} funzionante (16/20 arrivati: il percorso app$\to$frame$\to$RX$\to$app gira);
  \item \textbf{Inaffidabilit\`a} (4 pacchetti persi per bit-error sul link);
  \item \textbf{Nessuna ritrasmissione} (i seq persi non riappaiono mai --- con gli stream sarebbero stati
    ritrasmessi);
  \item \textbf{Nessun head-of-line blocking} (dopo il buco su 2, il seq 3 arriva immediatamente);
  \item \textbf{Negoziazione} corretta (con \code{max\_datagram\_frame\_size}=0 sarebbero stati scartati tutti;
    invece 16 sono passati $\Rightarrow$ il peer ha annunciato supporto).
\end{itemize}

% =====================================================================
\section{Build ed esecuzione}
Il workspace \`e gestito da \code{opp\_env} (ambiente Nix riproducibile): il build va fatto entro la sua shell.

\begin{lstlisting}[style=sh]
# build incrementale di INET (rigenera i makefile per i nuovi .cc, ~30s con ccache)
opp_env run -w . --no-build -c 'make -C inet-4.6.0 MODE=release -j 4'

# esecuzione del test funzionale
opp_env run -w . --no-build -c 'cd inet-4.6.0/examples/quic/datagram_test && \
  opp_run_release -u Cmdenv \
    -n <INET>/src:<INET>/examples \
    -l <INET>/out/clang-release/src/libINET.so omnetpp.ini -r 0'
\end{lstlisting}
\noindent Note: buildare dalla \emph{radice} di INET (non da \code{src/}) affinch\'e \code{make makefiles}
rigeneri la lista OBJS e includa i nuovi file. Nell'\code{omnetpp.ini} il livello di log va per-oggetto
(\code{**.cmdenv-log-level=info}).

% =====================================================================
\section{Versioning e prossimi passi}
Poich\'e INET \`e una dipendenza pinnata (non un repository git), tutte le modifiche sono conservate come
\textbf{patch versionata} nel repository \code{tod-omnet}: \code{tod-omnet/inet-patches/quic-datagram.patch}
(trasporto) e \code{tod-omnet/inet-patches/test/} (app di test e rete), con uno snapshot pristino in
\code{orig/} per rigenerare il diff.

\medskip
\noindent\textbf{Fase 2 (integrazione applicativa in tod-omnet).}
\begin{itemize}[nosep,leftmargin=1.4em]
  \item \code{TODCarApp} invia i dati-sensore via \code{sendDatagram} (mantenendo uno \code{streamId}
    ``nominale'' nel payload per identificare il sensore), spezzandoli in chunk $\le$ \code{max\_datagram\_frame\_size};
  \item un \emph{manifest} sullo stream~0 affidabile dichiara, per ogni sensore, numero di datagram e byte;
  \item \code{TODAgentApp} calcola un \code{loss\_ratio} \emph{pesato per-sensore};
  \item un interruttore di configurazione permette il confronto A/B tra ``datagram inaffidabile'' e ``stream
    affidabile''.
\end{itemize}

\end{document}
